\documentclass[12pt]{jlreq}
\usepackage{amsmath, amssymb}

\newcommand{\C}{\mathbb{C}}
\newcommand{\R}{\mathbb{R}}
\newcommand{\Z}{\mathbb{Z}}
\newcommand{\Tr}{\operatorname{Tr}}
\newcommand{\Hom}{\operatorname{Hom}}
\newcommand{\GL}{\operatorname{GL}}
\newcommand{\rank}{\operatorname{rank}}
\newcommand{\Id}{\operatorname{Id}}
\newcommand{\Ker}{\operatorname{Ker}}

\begin{document}

$\R^2 \times \C$ 上の複素数値関数
\[
  \varphi : ((x,y), \lambda) \in \R^2 \times \C \longrightarrow \varphi(x,y; \lambda) \in \C
\]
で，次の条件 (a), (b), (c) を同時にみたすものを求めよ。
\begin{enumerate}
  \item[(a)] $p, q$ を $p \ne q$ をみたす与えられた実数とし，
  \[
    \tau(x,y) = 1 + \exp\left[ (p-q)x + (p^2-q^2)y \right],\quad u(x,y) = \frac{\partial^2}{\partial x^2} \log\tau(x,y)
  \]
  とすると，$\varphi$ は平面 $\R^2$ で定義された偏微分方程式
  \[
    \frac{\partial\varphi(x,y;\lambda)}{\partial y} = \frac{\partial^2\varphi(x,y;\lambda)}{\partial x^2} + 2u(x,y)\varphi(x,y;\lambda) + 2\lambda \frac{\partial\varphi(x,y;\lambda)}{\partial x}
  \]
  の解になる。
  \item[(b)] 任意の $(x,y) \in \R^2$ に対して，
  \[
    \lim_{\lambda\to\infty} \varphi(x,y; \lambda) = 1
  \]
  \item[(c)] $(x,y) \in \R^2$ を任意に固定すると，$\varphi$ は $\lambda$ に関する有理関数であって，複素平面 $\C$ 上の特異点は一位の極ひとつのみであり，その極の位置は $(x,y)$ に依存しない。
\end{enumerate}

\end{document}